\subsection{RPLIDAR A1M8}

\vspace{2em}

\subsubsection{Description et fonctionnement}

Le RPLIDAR A1M8 est un capteur LiDAR 2D capable de scanner à 360$^\circ$ dans un plan horizontal. Il utilise un laser infrarouge de classe 1 (sans danger pour l’œil) et fonctionne par triangulation optique : il envoie un faisceau laser et calcule la distance à l’obstacle en analysant l’angle de retour.

\vspace{1em}

Il fournit environ 1450 mesures par tour, avec une résolution proche de 1$^\circ$ et une fréquence de scan autour de 5{,}5~Hz. On peut aussi augmenter la vitesse de rotation (jusqu’à 10~Hz) en réduisant le nombre de points. Sa portée utile est d’environ 6~mètres, avec un maximum pouvant atteindre 12~mètres selon les versions.

\vspace{1em}

Les données (angle + distance) permettent de créer une carte 2D de l’environnement, utile pour la détection d’obstacles, le suivi de mouvement ou des applications type SLAM. Il se connecte facilement à une Raspberry Pi via USB et fonctionne en plug-and-play. Grâce à la technologie OPTMAG (transmission sans fil entre les parties fixes et mobiles), il évite les frottements mécaniques, ce qui améliore la durabilité.

\vspace{1em}

Le capteur fonctionne aussi dans le noir, car il génère sa propre lumière, mais reste sensible aux fortes sources lumineuses parasites.

\vspace{2em}

\subsubsection{Avantages et inconvénients}

Le RPLIDAR A1M8 présente plusieurs atouts intéressants pour la détection de l’environnement. Il permet un balayage complet à 360$^\circ$, ce qui évite les angles morts et offre une bonne perception de l’espace autour de l’utilisateur. Il produit un nuage de points dense et précis, avec une résolution fine et une portée correcte pour un capteur abordable. Grâce à son taux d’échantillonnage élevé (jusqu’à 8000 mesures par seconde), il peut suivre les changements de l’environnement en temps réel. Contrairement à une caméra, il fournit directement des distances fiables, indépendamment de la lumière ambiante, ce qui le rend utilisable dans l’obscurité.

\vspace{1em}

Dans notre projet, il permet de détecter les obstacles proches et d’améliorer la précision locale, là où le GPS est trop imprécis ou l’IMU sujette à la dérive.

\vspace{1em}

Cependant, ce capteur a aussi ses limites. Il est 2D: il ne capte que ce qui se trouve dans son plan de balayage horizontal, ce qui peut lui faire rater des objets trop bas ou trop hauts. Sa portée est limitée à environ 6~mètres, ce qui restreint son usage à des environnements proches. Sa résolution angulaire (environ 1$^\circ$) peut laisser passer de petits objets à distance, et il est sensible à la nature des surfaces : les objets sombres, brillants ou transparents peuvent ne pas être bien détectés.

\vspace{1em}

Les conditions météo comme la pluie ou le brouillard réduisent aussi la qualité des mesures. En plein soleil, le capteur peut être perturbé par l’infrarouge naturel, ce qui crée du bruit et fausse les résultats. Enfin, étant un dispositif rotatif, il consomme plus d’énergie, produit beaucoup de données à traiter en continu, et subit une usure mécanique à long terme.

\vspace{1em}

Malgré tout, son bon rapport qualité/prix et sa compatibilité avec la Raspberry Pi en font un outil efficace, à condition de bien gérer ses limites techniques.

\vspace{2em}

\subsubsection{Sources de bruitage et défauts techniques}

Le RPLIDAR A1M8 peut être perturbé par plusieurs facteurs. D’abord, il y a un léger bruit aléatoire sur les distances mesurées, même sur une cible fixe. Plus la distance augmente, plus l’écart entre les mesures peut devenir important. Ensuite, les objets sombres ou peu réfléchissants sont souvent mal détectés : par exemple, un tissu noir peut parfois ne pas apparaître du tout dans le scan, contrairement à une surface claire au même endroit.

\vspace{1em}

Le capteur a aussi du mal avec les surfaces inclinées ou les objets très proches quand l’angle du faisceau est trop rasant. De plus, une lumière solaire directe peut saturer le capteur, entraînant des mesures fausses ou des pertes de détection. Ce phénomène a été observé lorsque le LiDAR était orienté vers le soleil.

\vspace{1em}

Enfin, comme le capteur est mécanique, de petites vibrations ou des mouvements pendant le scan peuvent légèrement déformer le nuage de points. Ces limites nous ont poussés à calibrer le capteur et à appliquer un filtrage pour améliorer la fiabilité des données.

\vspace{2em}

\vspace{1em}

\textbf{Remarque :} Bien que le capteur RPLIDAR A1M8 fasse partie du matériel fourni dans le cadre du projet, nous n’avons finalement pas intégré son utilisation dans notre système. Nous avons choisi de concentrer nos efforts sur l’exploitation combinée du GPS et de l’IMU, qui suffisent à eux seuls à réaliser un suivi de trajectoire pertinent pour notre cas d’usage. Par conséquent, aucun travail de calibrage ou d’intégration logicielle du LiDAR n’a été réalisé dans notre pipeline. Ce choix nous a permis de focaliser le développement sur l’amélioration de la précision de fusion GPS/IMU, tout en assurant un système fonctionnel et maîtrisé.

